Building an auditable mapping requires first a principled organization of the evaluation literature. We group technical AI evaluation methods into twelve families by \emph{what they measure}. Table~\ref{tab:taxonomy} summarizes the families; the subsections below describe each, its representative methods, its state of the art for LLM and agentic systems, and its maturity. The families form the rows of the mapping matrix in Section~\ref{sec:mapping}.

% T3 — taxonomy summary
\begin{table*}[!tp]
\centering
\caption{Taxonomy of technical AI evaluation method families. Maturity is assessed for the
LLM/agentic setting; classical-ML maturity is higher in several families. \textbf{Takeaway:}
most families are mature for classical/static models but only \emph{emerging} for LLM and
agentic systems---the maturity frontier, not the method inventory, is where conformity evidence
is currently thin.}
\label{tab:taxonomy}
\footnotesize
\begin{tabular}{p{0.04\textwidth}p{0.17\textwidth}p{0.40\textwidth}p{0.20\textwidth}}
\toprule
ID & Family & Representative methods / metrics & Maturity (LLM/agentic)\\
\midrule
M1 & Performance \& benchmarking & Held-out/OOD test sets, AUROC/F1, significance testing; MMLU, SuperGLUE, BIG-bench, agentic task suites & Mature (classical); evolving (agentic)\\
M2 & Fairness \& bias & Group/individual/counterfactual fairness, disparate impact; BBQ, StereoSet; action-level disparity & Mature (model); emerging (agentic)\\
M3 & Robustness \& distribution shift & PGD/C\&W attacks, randomized smoothing, ImageNet-C, WILDS, prompt-robustness & Mature (vision); emerging (LLM)\\
M4 & Safety evaluation & Hazard analysis, dangerous-capability evals, content-harm benchmarks, safety cases & Emerging\\
M5 & Red-teaming \& adversarial & Structured/automated red-teaming, jailbreaks, (indirect) prompt injection, agent tool-misuse & Emerging, fast-moving\\
M6 & Interpretability \& explainability & Feature attribution (SHAP, LIME, integrated gradients), concept testing (TCAV), probing, mechanistic interp., explanation-faithfulness, CoT faithfulness & Maturing; faithfulness contested\\
M7 & Uncertainty \& calibration & ECE/Brier, conformal prediction, selective prediction, semantic entropy, hallucination detection & Mature (calibration); emerging (LLM)\\
M8 & Data quality, governance \& docs & Datasheets, data cards, provenance audits, contamination/leakage detection, representativeness & Maturing\\
M9 & Transparency \& model documentation & Model/system cards, FactSheets, transparency indices, reproducibility checklists & Maturing\\
M10 & Human oversight \& interaction & Oversight-effectiveness, automation-bias and appropriate-reliance studies, escalation/deferral & Emerging\\
M11 & Post-deployment monitoring \& assurance & Drift detection, runtime monitoring, logging, incident reporting, third-party audit & Emerging\\
M12 & Privacy evaluation & Membership-inference and extraction attacks, differential-privacy verification, model-inversion resistance, privacy-risk assessment & Mature (attacks); emerging (LLM/agentic)\\
\bottomrule
\end{tabular}
\end{table*}




\subsection{M1: Performance \& benchmarking evaluation}

\paragraph{What it measures and why it matters.} The performance and benchmarking family quantifies how accurately and reliably a system maps inputs to correct outputs on held-out and out-of-distribution (OOD) data. This is the most basic conformity claim a deployer can make: that the system performs as intended on representative data and degrades gracefully under shift. Standardized metrics with appropriate statistical treatment are the operational counterpart to the accuracy and documentation expectations of the NIST AI RMF \cite{nist2023airmf}, the EU AI Act \cite{euaiact2024}, and ISO/IEC 42001 \cite{iso42001}.

\paragraph{Representative methods, metrics, and tools.} Classical supervised evaluation reports threshold-based and ranking metrics (accuracy, precision/recall, F1, AUROC) with significance testing on held-out splits, and increasingly assesses OOD generalization and adversarial robustness under shift \cite{li2024oodobustbench}. Calibration is examined alongside raw accuracy to ensure reported confidences are trustworthy \cite{le2024revisiting}. Provenance and reproducibility are documented through model cards \cite{mitchell2019modelcards} and datasheets \cite{gebru2018datasheets}, which record evaluation conditions and intended use. For language systems, aggregate competence is measured on multi-task suites such as MMLU \cite{hendrycks2020mmlu}, SuperGLUE \cite{wang2019superglue}, and the broad BIG-bench collection \cite{srivastava2023bigbench}.

\paragraph{LLM and agentic state of the art.} Contemporary practice extends benchmarking to generative and agentic behavior: instruction-following is scored with debiased automatic evaluators \cite{dubois2024alpacaeval}, factuality with hard multi-turn hallucination benchmarks \cite{hendrycks2024haulaeval}, and prompt-level robustness with adversarial perturbation suites \cite{dong2023promptrobust}. Agentic systems are assessed on end-to-end task completion across architectures \cite{anil2024agentarch}, multi-dimensional enterprise criteria \cite{song2025multidimensional}, and multilingual performance-and-security settings \cite{wang2023maps}. Structured high-stakes evaluation protocols (including a recent unpublished framework; citation withheld for double-blind review) organize these heterogeneous benchmarks into auditable evidence.

\paragraph{Maturity.} Classical ML benchmarking is mature, whereas LLM and agentic benchmarking remains rapidly evolving and contested, with active debate over contamination and construct validity. As summarized in Table~\ref{tab:matrix}, this family principally supplies evidence for accuracy, robustness, and performance-monitoring requirements across all three regimes.

\subsection{M2: Fairness \& bias evaluation}

This family measures whether a system's outputs, predictions, or actions systematically advantage or disadvantage groups defined by protected attributes. Anti-discrimination is an explicit demand across all three regimes: fairness and harmful-bias management is a core NIST trustworthiness characteristic \cite{nist2023airmf}, the EU AI Act requires examination for discriminatory outcomes and appropriate data governance \cite{euaiact2024}, and ISO/IEC 42001 obliges organizations to treat fairness-related impacts within the management system \cite{iso42001}.

\paragraph{Representative methods and metrics.} Classical evaluation distinguishes group fairness, individual fairness, and causal notions, surveyed comprehensively by \textit{Mehrabi et al.} \cite{mehrabi2021survey}. Group criteria include equalized odds and equality of opportunity \cite{hardt2016equality}, while individual fairness formalizes the principle that similar individuals receive similar treatment \cite{dwork2012fairness}. Counterfactual fairness instead conditions on causal interventions over protected attributes \cite{kusner2017counterfactual}. A well-known result is that several group criteria are mutually incompatible except in degenerate cases \cite{kleinberg2017inherent,chouldechova2017fair}, so conformity evidence must state which definition is being satisfied. Foundational audits such as \textit{Gender Shades} \cite{buolamwini2018gender} and embedding-bias analyses \cite{caliskan2017semantics,bolukbasi2016man} established disparate-impact measurement in practice.

\paragraph{LLM and agentic state of the art.} For language models, fairness evaluation has shifted toward benchmark suites: BBQ \cite{parrish2022bbq}, StereoSet \cite{nadeem2021stereoset}, CrowS-Pairs \cite{nangia2020crows}, and BOLD \cite{dhamala2021bold}, consolidated in recent surveys \cite{gallegos2024bias,chu2024fairness,navigli2023biases} and compositional or trust-oriented batteries \cite{wang2024ceb,huang2024trustllm,wang2023decodingtrust}, alongside dialect-level disparities in reward models \cite{mire2025rejected}. The agentic frontier moves from text to decisions: action-level disparity \cite{li2025actions}, tool-selection bias \cite{blankenstein2025biasbusters}, multi-agent fairness frameworks \cite{ranjan2025fairness}, and a recent agent-fairness framework (citation withheld for double-blind review), which attributes per-component fairness across an agent pipeline.

\paragraph{Maturity.} Static-classifier fairness is mature with established metrics, whereas LLM and agent action-level fairness remain emerging and benchmark-sensitive. This evidence speaks primarily to the discrimination and fairness obligations mapped in Table~\ref{tab:matrix} across NIST, the EU AI Act, and ISO/IEC 42001.

\subsection{M3: Robustness \& distribution-shift evaluation}

This family measures whether a model's predictive quality and safety properties persist when inputs deviate from the training distribution, whether through naturally occurring shift, accidental corruption, or adversarial manipulation. Because high-risk systems are rarely deployed in conditions identical to those under which they were trained, robustness evidence directly substantiates claims about reliability and resilience under foreseeable misuse and operating conditions \cite{nist2023airmf,euaiact2024}.

\paragraph{Representative methods and metrics.} Adversarial robustness has a mature evaluative core: gradient-based attacks such as PGD provide a standard adversarial-accuracy stress test \cite{madry2018pgd,goodfellow2015explaining}, and stronger optimization-based attacks expose models whose apparent robustness is an artifact of weak attacks \cite{carlini2016c2}. Certified methods complement empirical attacks by yielding provable radii within which predictions cannot change, most prominently randomized smoothing \cite{cohen2019certified}. Corruption robustness is quantified against systematic perturbation suites \cite{hendrycks2019corruption} and naturally occurring hard cases \cite{hendrycks2019natural}, while distribution-shift generalization is assessed on curated in-the-wild benchmarks \cite{koh2021wilds} and through methods that target invariance across environments \cite{arjovsky2019irm}. Out-of-distribution detection \cite{liang2024ood} and test-time adaptation, evaluated under realistic unsupervised hyperparameter selection \cite{tta2024realistic}, address whether a system can flag or absorb inputs it should not confidently process.

\paragraph{LLM and agentic state of the art.} For language models, robustness evaluation has shifted toward prompt sensitivity and jailbreak resistance. Standardized jailbreak benchmarks \cite{jailbreakbench2024} and randomized-perturbation defenses \cite{smoothllm2023} measure resilience to adversarial prompting, prompt-sensitivity probes quantify brittleness under benign paraphrase \cite{brittlebench2025}, and data-poisoning benchmarks assess training-time vulnerability \cite{poisonbench2024}; the generative-AI profile frames these as managed risks \cite{nist2024genai}.

\paragraph{Maturity.} Vision robustness and certified methods are mature, whereas LLM and agentic prompt-robustness evaluation remains emerging. As mapped in Table~\ref{tab:matrix}, this evidence speaks most directly to accuracy, robustness, and cybersecurity obligations \cite{euaiact2024}, the NIST resilience and risk-management functions \cite{nist2023airmf}, and operational controls \cite{iso42001}.

\subsection{M4: Safety evaluation}

\paragraph{What it measures and why it matters} The safety family quantifies a system's propensity to cause or enable harm, spanning hazard identification, dangerous-capability elicitation, content-harm measurement, and the structured argumentation that ties such evidence to an acceptable-risk conclusion. It underpins the risk-management and harm-mitigation obligations of all three regimes \cite{nist2023airmf,euaiact2024,iso42001}, and for frontier models is the principal artifact showing that residual hazards have been characterized rather than merely asserted.

\paragraph{Representative methods, metrics, and tools} Content-harm benchmarks operationalize safety as measurable rates of unsafe responses: \textit{SafetyBench} structures multiple-choice probes across harm categories \cite{wang2023safetybench}, while \textit{TruthfulQA} targets the related failure of confidently reproducing human falsehoods \cite{lin2021truthfulqa}. Mitigation-oriented training such as Safe RLHF separates helpfulness from harmlessness so that safety can be tuned and audited as an explicit objective \cite{dai2024saferlhf}. At the assurance level, structured safety cases assemble claims, arguments, and evidence into an auditable conformity narrative \cite{hawkins2025safetycaseai}, and risk-perspective analyses examine whether layered alignment mechanisms constitute independent defenses or share common failure modes \cite{gal2025alignment}.

\paragraph{LLM and agentic state of the art} Frontier evaluation has shifted toward eliciting \textit{dangerous capabilities} (for example cyber-offense, deception, and self-proliferation) under adversarial conditions to bound worst-case misuse potential \cite{hendel2024dangerouseval}. Because such measurements are sensitive to attack strategy, jailbreak-evaluation toolkits standardize how successful circumvention is judged and aggregated \cite{akhtar2024jailbreakeval}, and the NIST Generative AI Profile catalogs the corresponding risk categories and suggested actions \cite{nist2024genai}. For agentic systems, capability elicitation increasingly targets multi-step tool-use trajectories rather than single completions.

\paragraph{Maturity} Content-harm benchmarking and red-team-style jailbreak evaluation are comparatively mature, whereas dangerous-capability elicitation and quantified safety cases remain emerging and contested. As mapped in Table~\ref{tab:matrix}, this evidence speaks most directly to risk-identification, harm-mitigation, and accuracy--robustness obligations across the three frameworks \cite{nist2023airmf,euaiact2024,iso42001}.

\subsection{M5: Red-teaming \& adversarial testing}

This family measures whether a system can be induced to violate its intended safety policy under adversarial pressure, rather than under benign use. It probes for failure modes that ordinary functional testing does not surface: jailbreaks that bypass alignment guardrails, prompt injection (including \textit{indirect} injection routed through retrieved or tool-returned content), and tool-misuse by autonomous agents. Regulators frame robustness and security against adversarial manipulation as an obligation distinct from accuracy, to be assessed before and during deployment \cite{euaiact2024,nist2023airmf}.

\paragraph{Representative methods, metrics, and tools.} Optimization-based attacks generate transferable adversarial suffixes that defeat aligned models \cite{zou2023universal}, while query-efficient black-box methods elicit harmful completions in few interactions \cite{carlini2023extracting}. Standardized harnesses make these results comparable: HarmBench formalizes automated red-teaming and robust-refusal scoring across attack and defense pairs \cite{mazeika2024harmbench}, and benchmark suites quantify instruction-following robustness to prompt injection \cite{hendrycks2023advbench}. Indirect injection has been characterized both as a real-world compromise vector for LLM-integrated applications \cite{greshake2023notwhat} and through dedicated benchmarks that test whether firewalls or stronger evaluations are needed \cite{wang2024indirect}, with prompt-tuning shields proposed as lightweight countermeasures \cite{rai2024soft}. Typical metrics are attack success rate, refusal rate, and transferability across models and prompts.

\paragraph{LLM and agentic state of the art.} Red-teaming is increasingly automated and learning-based, with adversarial generators trained to discover robustness failures at scale \cite{kang2025learning,liu2025automatic}. Agentic evaluation extends the threat surface to tool calls and actions: Agent-SafetyBench measures unsafe agent behavior \cite{zhang2024agent}, SafeSearch red-teams search agents \cite{gu2025safesearch}, and EU-Agent-Bench scores illegal agent actions against legal baselines \cite{peng2025eu}. Recent systematization consolidates these methodologies for assurance and security \cite{wang2026systematic}.

\paragraph{Maturity.} Attack methods and benchmarks are mature for single-turn LLMs but emerging for indirect-injection and agentic tool-misuse settings, where standardization is incomplete. As mapped in Table~\ref{tab:matrix}, it evidences robustness and adversarial-security requirements across the three frameworks and their operational controls \cite{nist2024genai,iso42001}.

\subsection{M6: Interpretability \& explainability evaluation}

This family measures the extent to which a model's internal computations and outputs can be explained, attributed to inputs, and verified as faithful to the model's actual decision process. Such evidence underpins transparency, human oversight, and the ability to interrogate a decision; it speaks to the explainability and transparency expectations of the NIST AI RMF \cite{nist2023airmf}, the documentation and human-oversight obligations of the EU AI Act \cite{euaiact2024}, and the operational-control evidence required by ISO/IEC 42001 \cite{iso42001}.

\paragraph{Representative methods and metrics.} Feature-attribution techniques such as integrated gradients provide axiomatically grounded importance scores \cite{sudre2024axiomatic_integrated_gradients}, while concept-based approaches like Testing with Concept Activation Vectors (TCAV) quantify a model's sensitivity to human-interpretable concepts beyond raw features \cite{kim2019tcav}. Probing classifiers assess what linguistic and structural information is encoded in contextual representations \cite{tenney2019edge_probing}, building on the transformer architecture that motivates much of this analysis \cite{vaswani2017transformer}. A recurrent concern is whether such explanations are faithful rather than merely plausible: attention weights, for instance, have been shown not to constitute reliable explanations \cite{jain2019attention_not_explanation}. Benchmarking efforts such as XAI-Units introduce unit-test-style evaluation to verify that explainability methods behave as specified \cite{yang2024xai_units}.

\paragraph{LLM and agentic state of the art.} For generative and reasoning models, evaluation has shifted toward the faithfulness of chain-of-thought (CoT) traces \cite{yang2023measuring_faithfulness_cot}, with work distinguishing faithfulness from surface plausibility \cite{turpin2024faithfulness_plausibility} and arguing that self-explanations can still help predict model behavior \cite{lanham2024positive_case_faithfulness}. Mechanistic interpretability increasingly probes the internal circuitry of CoT via sparse autoencoding \cite{zou2024cot_mechanistic} and applies causal tracing to localize object representations and mitigate hallucination in vision-language models \cite{zhang2024causal_tracing_vqa}.

\paragraph{Maturity.} Attribution and probing are relatively \textit{mature}, whereas CoT-faithfulness and mechanistic interpretability remain \textit{emerging}. As mapped in Table~\ref{tab:matrix}, this family supplies evidence for transparency, explainability, and human-oversight requirements across all three regimes \cite{nist2023airmf,euaiact2024,iso42001}.

\subsection{M7: Uncertainty quantification \& calibration}

This family measures whether a model's stated confidence corresponds to its empirical correctness, and whether it can abstain or flag low-confidence outputs rather than asserting them. Regulators expect high-risk systems to expose the limits of their reliability: NIST frames calibrated uncertainty under the \textit{Measure} function as a precondition for validly characterized systems \cite{nist2023airmf}, the EU AI Act ties accuracy and robustness disclosures to technical documentation and instructions of use \cite{euaiact2024}, and ISO/IEC 42001 anchors them in performance evaluation and monitoring \cite{iso42001}.

\paragraph{Representative methods and metrics.} Classical calibration is assessed through expected calibration error (ECE) and Brier-type scores, with temperature scaling established as a strong post-hoc remedy for the overconfidence of modern networks \cite{guo2017calibration}, and broader surveys cataloguing the state of the art \cite{wang2023calibration}. Predictive uncertainty itself is estimated via Bayesian approximations such as Monte Carlo dropout \cite{gal2016bayesian}, evidential formulations for classification \cite{sensoy2018evidential} and regression \cite{amini2020deepevidential}, and learned confidence for out-of-distribution detection \cite{devries2018confidence}. Conformal prediction supplies distribution-free coverage guarantees and abstention-ready prediction sets \cite{angelopoulos2023conformal}, with work disentangling its interaction with deep model uncertainty \cite{barella2023quantifying} and isolating epistemic components \cite{ravi2026epistemic}; such sets have been shown to improve human decision making \cite{cresswell2024conformal}.

\paragraph{LLM and agentic state of the art.} For generative models, semantic uncertainty marginalizes over meaning-equivalent generations \cite{kuhn2023semantic}, with cheaper semantic-entropy probes \cite{kossen2024semantic} and token-level temperature scaling refining the estimate \cite{lamb2026temperature}. Zero-resource consistency checks detect hallucination without ground truth \cite{manakul2023selfcheckgpt}, while selective prediction reduces unnecessary abstention in multimodal reasoning \cite{srinivasan2024selective}; the GenAI Profile elevates confabulation and factuality as named risks \cite{nist2024genai}.

\paragraph{Maturity.} Tabular and vision calibration and conformal methods are mature, whereas LLM-level semantic uncertainty and hallucination detection remain emerging and contested. As mapped in Table~\ref{tab:matrix}, this evidence speaks chiefly to accuracy, robustness, and reliability-disclosure obligations across all three frameworks.

\subsection{M8: Data quality, governance \& documentation}
\label{subsec:m8}

This family assesses the data on which high-risk AI systems are built and the artifacts that record how that data was collected, curated, and used. Rather than probing model behavior directly, it measures upstream properties: quality dimensions such as completeness, accuracy, and consistency; representativeness relative to the deployment population; provenance and licensing of training corpora; and the existence and quality of documentation. Regulators treat data governance as a precondition for trustworthy outputs: the EU AI Act requires training, validation, and testing datasets to be relevant, sufficiently representative, and examined for bias \cite{euaiact2024}; the NIST AI RMF locates these concerns in its \textit{Map} and \textit{Measure} functions \cite{nist2023airmf}; and ISO/IEC 42001 requires documented controls over data resources and their lifecycle \cite{iso42001}.

\paragraph{Representative methods and tools.} Structured documentation artifacts are the most established evidence type, spanning Data Cards \cite{pushkarna2022datacards}, model and audit cards \cite{liang2024whatsdocumented,staufer2025auditcards}, and broader transparency instruments \cite{raffel2024transparency,winecoff2024governance}. Data-quality assessment is supported by surveys of quality dimensions and tooling \cite{pan2024metrics}, representativeness and downstream-fairness analyses \cite{liang2024representativeness}, and large-scale provenance and licensing audits \cite{luccioni2023dataprovenance}; provenance can also be captured automatically at the framework level \cite{pocock2021tribuo}. Benchmark and dataset catalogues further inform what evaluation data is available and appropriate \cite{ivanov2024benchmarks}.

\paragraph{LLM and agentic state of the art.} For large language models, the dominant emerging concern is data contamination, where benchmark items leak into pretraining corpora and inflate measured performance; detection methods and their limits are now actively surveyed \cite{lukoševičius2024contamination}. The opacity of web-scale corpora also makes provenance, licensing, and representativeness substantially harder to establish than in the curated-dataset setting \cite{luccioni2023dataprovenance,nist2024genai}.

\paragraph{Maturity.} Documentation templates and quality-dimension tooling are mature, whereas contamination detection, automated provenance tracking, and quantitative representativeness auditing remain emerging \cite{global2024governance}. As mapped in Table~\ref{tab:matrix}, it speaks most directly to data-governance, dataset-representativeness, record-keeping, and transparency requirements.

\subsection{M9: Transparency \& model documentation}
\label{subsec:m9}

This family measures whether an AI system's design, training data, intended use, evaluation results, and known limitations are recorded in a structured, verifiable, and externally legible form. Unlike performance families that score model behavior, M9 assesses the quality and completeness of the \textit{evidentiary record} itself: documentation is the connective tissue that lets an auditor trace a claim back to the test that supports it. This makes the family foundational for conformity, because most regulatory regimes treat the absence of adequate documentation as a deficiency in its own right, independent of how the underlying system performs.

Representative methods center on standardized artifacts and the pipelines that produce and grade them. AI FactSheets propose a supplier-side methodology for declaring system facts across the lifecycle \cite{weidinger2024aifactsheets}, while EvalCards standardize how evaluation results are reported \cite{lin2025evalcards} and model-reporting proposals explicitly merge EU regulatory fields into developer documentation \cite{pan2024modelreporting}. Scoring instruments operationalize transparency as a measurable quantity: the Foundation Model Transparency Index ranks developers against disclosure indicators \cite{wan2025fmti}, and the AI Transparency Atlas adds a real-time model-card evaluation pipeline \cite{shah2025aitransparencyatlas}. Domain extensions, such as bias-aware clinical model cards \cite{tong2024benchmarking}, and global information-exchange frameworks \cite{weidinger2023framework} broaden the documentation surface, while reproducibility programs supply the experimental-evidence backbone \cite{pineau2020reproducibility}. A recent unpublished framework (citation withheld for double-blind review) contributes a structured documentation-and-reporting schema for high-stakes systems.

For LLMs and agentic systems the state of the art is consolidating: surveys map explainability and transparency techniques for large models \cite{zhang2025transparentai,liang2023aiglmsurvey}, empirical-study guidelines target LLM evidence quality \cite{mehdipour2025evaluation}, benchmark-of-benchmarks audits scrutinize the validity of reported scores \cite{qian2025benchmark2}, and the AI Agent Index documents technical and safety features of deployed agents \cite{li2025aiagentindex}. Case studies confirm that documentation methodologies can be evaluated empirically \cite{yang2024evaluating}.

The family is \textit{mature} for static models and \textit{emerging} for agentic and continuously updated systems. As mapped in Table~\ref{tab:matrix}, it evidences NIST \textsc{Govern}/\textsc{Map} documentation functions \cite{nist2023airmf,nist2024genai}, EU AI Act technical-documentation and transparency obligations \cite{euaiact2024}, and the documented-information and lifecycle controls of ISO/IEC 42001 \cite{iso42001}.

\subsection{M10: Human oversight \& human--AI interaction evaluation}
\label{subsec:m10}

This family measures whether a human retained in the loop can meaningfully understand, supervise, and correct an AI system in operation. Rather than scoring the model in isolation, it evaluates the joint human--AI decision unit along dimensions such as oversight effectiveness, automation bias, appropriate (calibrated) reliance, contestability, and the reliability of escalation or deferral pathways. Regulators treat human oversight as a primary risk control: it is a governance and measurement expectation under the NIST AI RMF \textit{Govern} and \textit{Measure} functions \cite{nist2023airmf}, a binding high-risk obligation under the EU AI Act \cite{euaiact2024}, and an operational control an AI management system must implement and audit under ISO/IEC 42001 \cite{iso42001}.

\paragraph{Representative methods, metrics, and tools.} Trust-calibration surveys frame the target construct as alignment between user trust and actual system reliability \cite{amershi2022trustcalibration}, and decision-theoretic formulations operationalize ``appropriate reliance'' as a measurable quantity rather than raw agreement \cite{guo2024reliance}. Reviews of automation bias catalogue when explanations help or harm oversight \cite{romeo2025automationbias}, motivating human-aligned XAI benchmarks \cite{wang2024xaibenchmark} and studies of how operators form mental models during collaboration \cite{wang2025mentalmodels}. Reliance drills probe over-dependence by injecting controllable errors \cite{chen2024reliancedrills}, while staged evaluation workflows integrate domain experts and lay users into the assessment loop \cite{liang2024stagedworkflows}, and system cards externalize oversight assumptions for review \cite{reddy2025systemcards}.

\paragraph{LLM and agentic state of the art.} For reasoning models, process-level supervision and step-by-step verification expose intermediate steps for human checking \cite{wang2023verifystepbystep,hou2025processrewardsurvey}, and interactive-alignment views treat oversight as continuous specification negotiation \cite{krause2023interactivealignment}. Agentic settings add escalation and deferral evaluation \cite{cui2024escalation}, action-level assessment beyond task completion \cite{tan2024agentassessment}, and constitution-based behavioral constraints \cite{tan2025constitutionalai}.

\paragraph{Maturity.} The family is \textit{emerging}: trust-calibration and automation-bias constructs are well established, but standardized oversight metrics for agentic systems remain immature. This evidence speaks primarily to the human-oversight and transparency requirements mapped in Table~\ref{tab:matrix}.

\subsection{M11: Post-deployment monitoring \& assurance}

This family measures whether a deployed system continues to behave within its validated operating envelope, and whether deviations are detected, logged, and escalated in a timely manner. Since pre-deployment evaluation captures only a snapshot, post-deployment monitoring is the evidentiary backbone of \textit{continuous} conformity: it speaks to the \textsc{Manage}-function operate-and-monitor obligations of the NIST AI RMF \cite{nist2023airmf}, to post-market monitoring, logging, and serious-incident reporting under the EU AI Act \cite{euaiact2024}, and to the performance-evaluation and continual-improvement clauses of ISO/IEC 42001 \cite{iso42001}.

\paragraph{Representative methods, metrics, and tools.} Core techniques include distribution-shift and concept-drift detection over inputs and learned representations \cite{greco2024driftlens,neuralnetwork2024robustness}, contextualized runtime monitoring that maps signals to operating context rather than raw statistics \cite{leest2025contextaware}, and uncertainty quantification used as a live trigger for shift and degradation \cite{neuralnetwork2024robustness}. Lifecycle-oriented frameworks organize these signals into structured assurance and traceable terminology \cite{aisystem2024evaluation}, while scalable policy-checking pipelines automate continuous compliance evaluation against codified requirements \cite{compliance2025pasta}. On the assurance side, third-party and frontier-model auditing methodologies \cite{auditing2024frontier}, layered governance-and-certification architectures \cite{audit2025responsible}, adaptive governance for decentralized deployment \cite{governance2025adaptive}, and incident-analysis taxonomies that convert reported harms into actionable patterns \cite{incidents2025harms} provide the external-evidence layer.

\paragraph{LLM and agentic state of the art.} For generative systems, monitoring increasingly targets factuality: decoding-time hallucination monitors that score partial responses during generation \cite{monitoring2025hallucination} and fact-checking pipelines that score factuality of deployed outputs \cite{mitigating2025hallucination}. For agentic systems, a recent agent-fairness framework (citation withheld for double-blind review) extends runtime monitoring to per-component, action-level disparity, attributing fairness drift to the responsible agent step rather than the aggregate output.

\paragraph{Maturity.} Drift detection and logging are mature; continuous agentic assurance, third-party audit, and live hallucination monitoring remain emerging. As summarized in Table~\ref{tab:matrix}, this evidence maps to the \textsc{Manage}/monitor, post-market monitoring and incident-reporting, and performance-evaluation requirements of the three regimes \cite{nist2023airmf,euaiact2024,iso42001}.

\subsection{M12: Privacy evaluation}

This family measures whether a system leaks information about the individuals in its
training data or the subjects of its inputs, and whether privacy-protective techniques
provide their claimed guarantees. Privacy is one of the trustworthiness characteristics
that high-risk regimes require be assessed: it is named explicitly in the NIST AI RMF
(\textsc{Measure}~2.10)~\cite{nist2023airmf}, underlies the data-governance and
transparency duties of the EU AI Act~\cite{euaiact2024}, and is a core impact dimension of
the ISO/IEC~42001 impact-assessment controls~\cite{iso42001}.

\paragraph{Representative methods and metrics.} The dominant empirical instruments are
attacks repurposed as audits. Membership-inference attacks estimate whether a given record
was in the training set and have become a standard privacy-leakage
metric~\cite{shokri2017membership}, with later work placing their evaluation on rigorous,
likelihood-ratio foundations to avoid overstated robustness~\cite{carlini2022firstprinciples}.
Training-data-extraction attacks demonstrate verbatim memorisation in large language
models, providing a direct, alarming measure of leakage~\cite{carlini2021extracting}, while
inference attacks show that models can expose attributes never explicitly
stated~\cite{staab2023memorization}. On the defensive side, differential privacy supplies a
provable guarantee whose privacy budget is itself the evaluation
metric~\cite{abadi2016dp}.

\paragraph{LLM and agentic state of the art.} Privacy evaluation for LLMs and agents is an
active and consolidating area, surveyed comprehensively in recent
work~\cite{neel2023privacy}; open problems include quantifying leakage through tool use and
retrieval, and evaluating memorisation at frontier scale.

A caveat governs how this family maps to conformity: most of its instruments are
\emph{vulnerability tests} (membership-inference, extraction, attribute-inference) that
\emph{detect} privacy leakage rather than \emph{assure} its absence. They evidence that a
privacy risk was examined, in keeping with the EU AI Act and ISO impact-assessment framing, but do not by
themselves demonstrate privacy adequacy; differential privacy is the exception, supplying a
provable guarantee. We therefore code privacy evaluation as having \emph{supporting} rather
than \emph{primary} relevance to the privacy requirement (NIST \textsc{Measure}~2.10).

\paragraph{Maturity.} Attack-based leakage measurement and differential-privacy accounting
are mature for classical models but unevenly applied to LLMs and agentic pipelines, and
this family is, relative to the others, comparatively thin in our corpus. As shown in
Table~\ref{tab:matrix}, it is the main \emph{partial}-evidence source for the privacy
requirements of all three frameworks; the privacy requirement itself is consequently
flagged as weakly served in the gap analysis (Section~\ref{sec:gaps}).
